
\documentclass{article}
\usepackage{colortbl}
\usepackage{makecell}
\usepackage{multirow}
\usepackage{supertabular}

\begin{document}

\newcounter{utterance}

\centering \large Interaction Transcript for game `hot\_air\_balloon', experiment `air\_balloon\_survival\_en\_negotiation\_easy', episode 4 with Qwen3{-}32B{-}t0.0.
\vspace{24pt}

{ \footnotesize  \setcounter{utterance}{1}
\setlength{\tabcolsep}{0pt}
\begin{supertabular}{c@{$\;$}|p{.15\linewidth}@{}p{.15\linewidth}p{.15\linewidth}p{.15\linewidth}p{.15\linewidth}p{.15\linewidth}}
    \# & $\;$A & \multicolumn{4}{c}{Game Master} & $\;\:$B\\
    \hline

    \theutterance \stepcounter{utterance}  
    & & \multicolumn{4}{p{0.6\linewidth}}{
        \cellcolor[rgb]{0.9,0.9,0.9}{
            \makecell[{{p{\linewidth}}}]{
                \texttt{\tiny{[P1$\langle$GM]}}
                \texttt{You are participating in a collaborative negotiation game.} \\
\\ 
\texttt{Together with another participant, you must agree on a single set of items that will be kept. Each of you has your own view of how much each item matters to you (importance). You do not know how the other participant values the items. Additionally, you are given the effort each item demands.} \\
\texttt{You may only agree on a set if the total effort of the selected items does not exceed a shared limit:} \\
\\ 
\texttt{LIMIT = 3145} \\
\\ 
\texttt{Here are the individual item effort values:} \\
\\ 
\texttt{Item effort = \{"C08": 948, "C55": 118, "A43": 469, "A67": 919, "B42": 284, "A62": 111, "B19": 806, "C43": 47, "C73": 848, "C10": 303, "B36": 13, "A18": 629, "C61": 687, "C67": 15, "C42": 94\}} \\
\\ 
\texttt{Here is your personal view on the importance of each item:} \\
\\ 
\texttt{Item importance values = \{"C08": 138, "C55": 583, "A43": 868, "A67": 822, "B42": 783, "A62": 65, "B19": 262, "C43": 121, "C73": 508, "C10": 780, "B36": 461, "A18": 484, "C61": 668, "C67": 389, "C42": 808\}} \\
\\ 
\texttt{Goal:} \\
\\ 
\texttt{Your goal is to negotiate a shared set of items that benefits you as much as possible (i.e., maximizes total importance to YOU), while staying within the LIMIT. You are not required to make a PROPOSAL in every message {-} you can simply negotiate as well. All tactics are allowed!} \\
\\ 
\texttt{Interaction Protocol:} \\
\\ 
\texttt{You may only use the following structured formats in a message:} \\
\\ 
\texttt{PROPOSAL: \{'A', 'B', 'C', …\}} \\
\texttt{Propose keeping exactly those items.} \\
\texttt{REFUSE: \{'A', 'B', 'C', …\}} \\
\texttt{Explicitly reject opponent's proposal.} \\
\texttt{ARGUMENT: \{'...'\}} \\
\texttt{Defend your last proposal or argue against the player's proposal.} \\
\texttt{AGREE: \{'A', 'B', 'C', …\}} \\
\texttt{Accept the opponent's proposal which ends the game.} \\
\texttt{STRATEGIC REASONING: \{'...'\}} \\
\texttt{	Describe your strategic reasoning or anticipation explaining your choice of action. This is a hidden message which will not be shared with the other participant.} \\
\\ 
\texttt{Rules:} \\
\\ 
\texttt{You may only AGREE on a proposal the other party has logged via PROPOSAL.} \\
\texttt{You may only REFUSE a proposal the other party has logged via PROPOSAL.} \\
\texttt{Total effort of any PROPOSAL or AGREE set must be ≤ LIMIT.} \\
\texttt{Do NOT reveal your hidden importance scores.} \\
\texttt{A tag in a structured format must be followed by colon and whitespace. The argument must be a python set containing 0 or more strings.} \\
\texttt{So, it must be of the form TAG: \{...\}} \\
\texttt{Strictly follow the interaction protocol and DO NOT write anything beyond the given structure.} \\
\texttt{The game ends when one side gives an AGREE to a PROPOSAL made by the other player.} \\
\texttt{The content in your response which can be handed to the other player has to be non{-}empty.} \\
\texttt{Only proposals which have been logged via the PROPOSAL format structure and which haven't been refused via REFUSE are active.} \\
\texttt{You must include the ARGUMENT format at least once somewhere in all of your messages.} \\
\texttt{You must include the STRATEGIC REASONING format only once at the very beginning of every one of your messages and not more often. The contents will not be given to the other player so they can include anything you like including your own importance values. Here you should reason step by step to come up with you next move.} \\
            }
        }
    }
    & \\ \\

    \theutterance \stepcounter{utterance}  
    & \multicolumn{4}{p{0.6\linewidth}}{
        \cellcolor[rgb]{0.9,0.9,0.9}{
            \makecell[{{p{\linewidth}}}]{
                \texttt{\tiny{[P1$\rangle$GM]}}
                \texttt{<think>} \\
\texttt{Okay, let's see. I need to figure out the best set of items to propose that maximizes my total importance without exceeding the effort limit of 3145. First, I should calculate the total effort for different combinations. Let me list out the items with their effort and importance values.} \\
\\ 
\texttt{Looking at the importance values, the highest ones are A43 (868), C55 (583), B42 (783), C10 (780), A67 (822), and C08 (138). Wait, A67 has high importance but also high effort (919). Let me check the effort for each. } \\
\\ 
\texttt{Let me start by considering the items with the highest importance. A43 has 868 importance and 469 effort. B42 is 783 importance and 284 effort. C10 is 780 importance and 303 effort. A67 is 822 importance but 919 effort. C55 is 583 importance and 118 effort. } \\
\\ 
\texttt{If I take A43, B42, and C10, their total effort would be 469 + 284 + 303 = 1056. That's under the limit. Adding A67 would add 919, which would make the total 1975, still under 3145. Wait, 1056 + 919 = 1975. Then maybe add C55 (118) to get 2093. Then maybe C42 (94) and C67 (15). Let's see: 2093 + 94 +15 = 2102. Still under. Then maybe B36 (13) and C43 (47). Adding those would be 2102 + 13 +47 = 2162. Still a lot left. } \\
\\ 
\texttt{But I need to maximize my total importance. Let me calculate the total importance for this combination. A43 (868) + B42 (783) + C10 (780) + A67 (822)} \\
            }
        }
    }
    & & \\ \\

    \theutterance \stepcounter{utterance}  
    & & \multicolumn{4}{p{0.6\linewidth}}{
        \cellcolor[rgb]{0.9,0.9,0.9}{
            \makecell[{{p{\linewidth}}}]{
                \texttt{\tiny{[P1$\langle$GM]}}
                \texttt{Your response did not start with the proper strategic reasoning tag at the very beginning of your response. The very first structured format must be of the form STRATEGIC REASONING: \{...\}. Try again.} \\
            }
        }
    }
    & \\ \\

    \theutterance \stepcounter{utterance}  
    & \multicolumn{4}{p{0.6\linewidth}}{
        \cellcolor[rgb]{0.9,0.9,0.9}{
            \makecell[{{p{\linewidth}}}]{
                \texttt{\tiny{[P1$\rangle$GM]}}
                \texttt{<think>} \\
\texttt{Okay, let me start over. The user wants me to make sure the first part is the STRATEGIC REASONING. So, I need to outline my thought process first.} \\
\\ 
\texttt{First, I need to maximize my total importance without exceeding the effort limit. Let me list the items by importance. The top ones are A43 (868), B42 (783), C10 (780), A67 (822), C55 (583), and C42 (808). Wait, C42 has high importance but low effort (94). } \\
\\ 
\texttt{I should prioritize items with high importance and low effort. Let's check their efforts. A43 is 469, B42 is 284, C10 is 303, A67 is 919 (high effort), C55 is 118, C42 is 94. } \\
\\ 
\texttt{If I take A43, B42, C10, and C42, their total effort is 469 + 284 + 303 + 94 = 1150. That's under 3145. Adding C55 (118) brings it to 1268. Then maybe add C67 (15) and B36 (13) for minimal effort. Total effort would be 1268 +15+13=1296. Still a lot left. } \\
\\ 
\texttt{But wait, A67 has high importance (822) but high effort (919). If I include it, the total effort would be 1150 +919=2069. Then adding C55 (118) makes 2187. Then C42 is already included. Maybe add C43 (47) and C61 (687). Wait, C61's effort is 687. Adding that would make 2187+687=2874. Still under 3145. } \\
\\ 
\texttt{But C61's importance is 668. So total importance would be A43(868) + B42(783) + C10(} \\
            }
        }
    }
    & & \\ \\

    \theutterance \stepcounter{utterance}  
    & & \multicolumn{4}{p{0.6\linewidth}}{
        \cellcolor[rgb]{0.9,0.9,0.9}{
            \makecell[{{p{\linewidth}}}]{
                \texttt{\tiny{[P1$\langle$GM]}}
                \texttt{Your response did not start with the proper strategic reasoning tag at the very beginning of your response. The very first structured format must be of the form STRATEGIC REASONING: \{...\}. Try again.} \\
            }
        }
    }
    & \\ \\

    \theutterance \stepcounter{utterance}  
    & \multicolumn{4}{p{0.6\linewidth}}{
        \cellcolor[rgb]{0.9,0.9,0.9}{
            \makecell[{{p{\linewidth}}}]{
                \texttt{\tiny{[P1$\rangle$GM]}}
                \texttt{<think>} \\
\texttt{Okay, I need to make sure I start with the STRATEGIC REASONING tag. Let me outline my plan step by step.} \\
\\ 
\texttt{First, I'll list all items with their importance and effort. The goal is to maximize my total importance without exceeding the effort limit of 3145. Let's sort the items by importance in descending order:} \\
\\ 
\texttt{1. A43: 868 (effort 469)} \\
\texttt{2. B42: 783 (284)} \\
\texttt{3. C10: 780 (303)} \\
\texttt{4. A67: 822 (919)} \\
\texttt{5. C42: 808 (94)} \\
\texttt{6. C55: 583 (118)} \\
\texttt{7. C61: 668 (687)} \\
\texttt{8. C08: 138 (948)} \\
\texttt{9. B36: 461 (13)} \\
\texttt{10. C67: 389 (15)} \\
\texttt{11. B19: 262 (806)} \\
\texttt{12. A18: 484 (629)} \\
\texttt{13. C43: 121 (47)} \\
\texttt{14. A62: 65 (111)} \\
\texttt{15. C73: 508 (848)} \\
\\ 
\texttt{Now, I need to select a combination that fits under 3145. Let's start with the highest importance items. A43, B42, C10, A67, C42, C55. Let's calculate their total effort:} \\
\\ 
\texttt{A43 (469) + B42 (284) = 753} \\
\texttt{+ C10 (303) = 1056} \\
\texttt{+ A67 (919) = 1975} \\
\texttt{+ C42 (94) = 2069} \\
\texttt{+ C55 (118) = 2187} \\
\\ 
\texttt{That's 2187. Remaining effort: 3145 {-} 2187 = 958. Can we add more? Let's see. Next highest is C61 (687). Adding that would bring total to} \\
            }
        }
    }
    & & \\ \\

    \theutterance \stepcounter{utterance}  
    & & & \multicolumn{2}{p{0.3\linewidth}}{
        \cellcolor[rgb]{0.9,0.9,0.9}{
            \makecell[{{p{\linewidth}}}]{
                \texttt{\tiny{[GM$|$GM]}}
                \texttt{abort} \\
            }
        }
    }
    & & \\ \\

    \theutterance \stepcounter{utterance}  
    & & & \multicolumn{2}{p{0.3\linewidth}}{
        \cellcolor[rgb]{0.9,0.9,0.9}{
            \makecell[{{p{\linewidth}}}]{
                \texttt{\tiny{[GM$|$GM]}}
                \texttt{end game} \\
            }
        }
    }
    & & \\ \\

\end{supertabular}
}

\end{document}
